\subsection{\fid{F62}: nonproper flag convolution and a nonconstant torus coefficient}
\label{sec:nonpropertorus}

The nonproper alternative can be tested with actual operations in
the standard rank-one flag geometry inside the ordinary Hecke
ambient setting. It realizes the raw factors in two specific
IV/VIII coefficients and reveals extension data invisible to
their pointwise traces. It does not identify the designated
nonstandard canonical labels or supply a nonstandard quadratic.

Fix \(k=\mathbb F_q\), a prime \(\ell\nmid q\), and a finite
\(\ell\)-adic coefficient field \(E\).
We use geometric Frobenius, so \(E(r)\) has eigenvalue \(q^{-r}\).
Put
\[
 X=\mathbb P^1_k,\qquad Y=X\times X,\qquad
 i:\Delta\simeq X\hookrightarrow Y,\qquad
 j:O=Y\setminus\Delta\hookrightarrow Y.
\]
The standard off-diagonal kernel is
\[
                         \mathcal H=j_!E_O.
\]
For constructible complexes on \(Y\), use convolution
\begin{equation}\label{eq:flagkernelconvolution}
 \mathcal A\star\mathcal B
   =Rp_{13!}(p_{12}^*\mathcal A\otimes^L p_{23}^*\mathcal B).
\end{equation}
Its unit is \(i_*E_X\).
These are actual geometric operations. Composition, proper base
change and the projection formula identify both parenthesizations
of three kernels with
\[
 Rp_{14!}(p_{12}^*\mathcal A\otimes
          p_{23}^*\mathcal B\otimes p_{34}^*\mathcal C).
\]
The functorial identifications for four factors give the
associativity pentagon. For powers of \(\mathcal H\), the common
correspondence is explicitly the variety of sequences
\((x_0,\ldots,x_r)\) with \(x_{a-1}\ne x_a\).
The required framework is proper base change, compact direct
images, and the projection formula
\cite[Theorem~17.10, pp.~116--117; Proposition~18.4,
pp.~121--122; Proposition~22.5, pp.~136--137]{Milne}.

For two factors define
\[
 Z=\{(x,y,z):x\ne y,\ y\ne z\}\subset X^3,\qquad
 f:Z\to Y,\quad f(x,y,z)=(x,z).
\]
Then \(\mathcal H\star\mathcal H=Rf_!E_Z\).
Over \(\Delta\) the fiber is \(\mathbb A^1\); over \(O\)
it is \(\mathbb G_m\).
Let
\[
 E_3=\{(x,y,z):x,y,z\text{ pairwise distinct}\},\qquad
 \pi:E_3\to O,\qquad \mathcal K=R\pi_!E.
\]

\begin{finding}[\fid{F62}, actual torus coefficient and nonsplitting]
\label{f:nonpropertorus}
There is an actual localization triangle
\begin{equation}\label{eq:flagopenclosedrelation}
 j_!\mathcal K\longrightarrow
 \mathcal H\star\mathcal H\longrightarrow
 i_*E_X(-1)[-2]\longrightarrow j_!\mathcal K[1].
\end{equation}
The off-diagonal family is the nontrivial torsor
\[
 E_3\simeq\mathcal L^\times,\qquad
                  \mathcal L=O(1,-1)|_O.
\]
Its coefficient complex fits in the triangle
\begin{equation}\label{eq:torusCherntriangle}
 \mathcal K\longrightarrow E_O(-1)[-2]
   \xrightarrow{c_1(\mathcal L)} E_O
   \longrightarrow\mathcal K[1],
\end{equation}
where \(c_1(\mathcal L)\ne0\).
In particular
\[
 \mathcal K\not\simeq E_O[-1]\oplus E_O(-1)[-2].
\]
Thus replacing the relative coefficient by the constant
complex of a \(\mathbb G_m\) fiber gives no direct-sum
convolution relation, despite agreement of all finite-extension
trace functions.
\end{finding}

\begin{proof}
\par\smallskip\noindent\textbf{The actual localization maps.}\quad
Restricting \(Rf_!E_Z\) to the two strata and applying the
open/closed localization triangle gives
\eqref{eq:flagopenclosedrelation}.
The diagonal family is an affine-line bundle; its only compact
direct image is the canonical top Tate line
\(E_X(-1)[-2]\).
This follows from compact-support base change and the trace
normalization for a smooth affine line
\cite[Proposition~18.4(b), pp.~121--122; Theorem~24.1(a),(b),
p.~144]{Milne}.
No product trivialization of the off-diagonal family is used.

For \(Q=q^n\), the trace of \(\mathcal H\) is the indicator
\(h\) of distinct rational points. Counting the intermediate
point gives
\begin{equation}\label{eq:unnormalizedflagHecke}
                         h*h=(Q-1)h+Q\,1_\Delta.
\end{equation}
Under \(v=Q^{1/2}\), the normalized operator \(T=v^{-1}h\)
satisfies
\[
 T^2=(v-v^{-1})T+1,\qquad (T-v)(T+v^{-1})=0.
\]
This is the standard ordinary Hecke normalization used in IV.
An arithmetic Weil line may implement \(v^{-1}\); no geometric
integral Tate object for a half twist is being asserted.
The unnormalized triangle needs no such coefficient choice.

\par\smallskip\noindent\textbf{Identifying the torsor and its orientation.}\quad
On \(O\), let \(L_x,L_z\) be the tautological lines from the
first and second factors. Then
\[
 \operatorname{Hom}(L_x,L_z)
      =\operatorname{pr}_1^*O(1)\otimes
        \operatorname{pr}_2^*O(-1)=\mathcal L.
\]
Two distinct lines \(x,z\subset k^2\) give \(k^2=x\oplus z\).
A third line \(y\) distinct from both is the graph of a unique
nonzero linear map \(x\to z\). Conversely such a graph is
distinct from both coordinate lines. These algebraic inverse
constructions commute with base change, proving
\(E_3\simeq\mathcal L^\times\).
This fixes \(\mathcal L=O(1,-1)\); the reverse graph convention
would give its dual and change the Chern-class sign.
Equivalently, this is the torsor
\(\mathrm{PGL}_2\to\mathrm{PGL}_2/T\), but the explicit graph
map supplies the identification.

Write \(h_x=c_1(O(1,0))\) and \(h_z=c_1(O(0,1))\).
The diagonal class is \(h_x+h_z\). Its Gysin sequence gives
\begin{equation}\label{eq:offdiagonalcohomology}
 H^0(O_{\bar k},E)=E,\qquad
 H^2(O_{\bar k},E)
   =\bigl(E(-1)h_x\oplus E(-1)h_z\bigr)/
                         E(-1)(h_x+h_z),
\end{equation}
with all other ordinary cohomology zero.
Indeed the first nonzero Gysin map sends \(1\) to
\(h_x+h_z\) and is injective; the map
\(H^2(\Delta,E)(-1)\to H^4(Y,E)\) is an isomorphism
because a point on the diagonal has degree one.
These use the standard smooth-pair and cycle-class normalization
\cite[Theorem~16.1 and Corollary~16.2, p.~108;
\(\S23\), pp.~138--139; Remark~24.2(b), p.~145]{Milne}.
Consequently
\begin{equation}\label{eq:torusnonzeroChern}
 c_1(\mathcal L)=h_x-h_z=2h_x\ne0
                       \quad\text{in }H^2(O_{\bar k},E(1)).
\end{equation}
Characteristic-zero coefficients make \(2\) nonzero even if
the finite field has characteristic two.

\par\smallskip\noindent\textbf{The connecting map.}\quad
Compactify the torsor to
\[
 p:\mathbb P(O_O\oplus\mathcal L)\to O
\]
with sections \(s_0=\mathbb P(O)\) and
\(s_\infty=\mathbb P(\mathcal L)\).
Their complement is \(\mathcal L^\times\).
Let \(\xi=c_1(O_{\mathbb P}(1))\), with projectivization
parametrizing lines. The unit and hyperplane maps give
\[
                   Rp_*E\simeq E_O\oplus E_O(-1)[-2].
\]
This is the projective-bundle formula, verified on the fibers
by proper base change. One can first apply it over the smooth
projective \(Y\), with \(O(1,-1)\), and then restrict to \(O\)
\cite[Theorem~23.2, pp.~139--140]{Milne}.

The restrictions of \(\xi\) to the two sections are \(0\) and
\(-c_1(\mathcal L)\), since \(O_{\mathbb P}(1)\) restricts
to \(O\) and \(\mathcal L^\vee\). Thus the restriction map to
the two sections has matrix
\begin{equation}\label{eq:twoSectionChernMatrix}
                    \begin{pmatrix}
                     1&0\\
                     1&-c_1(\mathcal L)
                    \end{pmatrix}.
\end{equation}
The Chern-class normalization and pullback are those of
\cite[Theorem~23.3(a),(b), p.~140]{Milne}.
Subtract the first row from the second, cancel the identity
summand, and change the sign of the remaining target.
The localization triangle becomes
\eqref{eq:torusCherntriangle}.
Changing that target identification back multiplies the
boundary by \(-1\) and gives an isomorphic triangle.
Thus the claimed nonzero boundary is an actual map, with
its harmless overall-sign choice specified.

It follows that
\[
 \mathcal H^1(\mathcal K)=E_O,\qquad
 \mathcal H^2(\mathcal K)=E_O(-1),
\]
and all other cohomology sheaves vanish. The nonzero
two-degree extension \eqref{eq:torusnonzeroChern} prevents
splitting.

For a direct check independent of extension terminology,
apply ordinary \(R\Gamma(O_{\bar k},-)\).
The map in degree two is the nonzero Chern map between
one-dimensional \(E(-1)\) spaces, hence an isomorphism.
The long exact sequence gives
\begin{equation}\label{eq:relativeTorusHypercohomology}
 H^1(O,\mathcal K)=E,\qquad
 H^4(O,\mathcal K)=E(-2),
\end{equation}
with other groups zero.
Its total dimension is two. For the constant split complex
\(E_O[-1]\oplus E_O(-1)[-2]\), the nonzero groups occur
in degrees \(1,2,3,4\), with types
\(E,E(-1),E(-1),E(-2)\), so the total dimension is four.
Here ordinary cohomology on the base of a relative compact
pushforward is not being identified with full compact
cohomology of \(E_3\).

In particular an isomorphism
\[
 \mathcal H\star\mathcal H\simeq
 \mathcal H\otimes\bigl(E[-1]\oplus E(-1)[-2]\bigr)
       \oplus i_*E_X(-1)[-2]
\]
would contradict the preceding nonsplitting after restriction
to \(O\).
Both sides nevertheless have the trace function
\eqref{eq:unnormalizedflagHecke} over every finite extension.

The outer triangle \eqref{eq:flagopenclosedrelation} is
itself nonsplit. Compactify \(Z\) in \(Y\times\mathbb P^1\).
Its complement is the closed union of the sections \(y=x\)
and \(y=z\), proper over \(Y\) with finite geometric fibers.
Its higher direct images vanish by proper base change.
The localization sequence therefore gives
\(R^2f_!E=E_Y(-1)\).
A split outer triangle would instead give
\[
 R^2f_!E=j_!E_O(-1)\oplus i_*E_X(-1).
\]
Both summands are nonzero, whereas the constant rank-one
sheaf on connected \(Y\) has endomorphism ring \(E\) and
is indecomposable. This contradiction proves nonsplitting
of the outer triangle as well as the separate coefficient
nonsplitting detected by the Chern class.
The valid geometric relation retains those extension maps.
\end{proof}

\paragraph{The actual IV and VIII scalar factors.}
Take two independent copies of the torsor, corresponding to
the two ordinary factors of the product ambient setting.
On \(O\times O\), put
\[
                   \mathcal K_2=(\mathcal K\boxtimes\mathcal K)(1).
\]
Every rational fiber over \(\mathbb F_{q^n}\) has trace
\begin{equation}\label{eq:relativeczerotrace}
 Q^{-1}(Q-1)^2=Q+Q^{-1}-2
             =(v-v^{-1})^2,\qquad v=Q^{1/2}.
\end{equation}
This is exactly IV's \(c_0=D-4\), the Figure~21 coefficient
in \fid{F58} \cite[Example~19.3, pp.~133--135]{IV}.
It is an actual relative compact pushforward of the product
torsor. The two localization triangles supply its boundary
maps; derived tensor product and its Koszul signs supply
their compatibility.

It is still not a constant coefficient complex.
Restrict to \(O\times\{o\}\) for a rational point \(o\).
After forgetting Weil structure, the fiber is
\(C=E[-1]\oplus E(-1)[-2]\).
The restriction is \(\mathcal K\otimes C\), up to twist;
its base hypercohomology has total dimension \(2\cdot2=4\)
by \eqref{eq:relativeTorusHypercohomology}.
A constant complex with the same stalk \(C\otimes C\)
has dimension \(2\cdot4=8\) on \(O\), by
\eqref{eq:offdiagonalcohomology}.

For the exact published coefficient in VIII,
\[
 c(\gamma_2^1,\gamma_1^1;\gamma_{21}^1)
   =-\frac{(q_{\rm src}^2+q_{\rm src}+1)(q_{\rm src}-1)^2}
           {q_{\rm src}^2},
\]
where \(q_{\rm src}\) is the deformation variable,
take the point complex
\begin{equation}\label{eq:VIIIgeometriccoefficient}
 \mathcal C_{\rm VIII}
   =R\Gamma_c((\mathbb G_m)^2\times\mathbb P^2,E(2))[1].
\end{equation}
Its trace, at \(q_{\rm src}=Q=q^n\), is exactly
\[
 -Q^{-2}(Q-1)^2(1+Q+Q^2)
       =-(v-v^{-1})^2(v^2+1+v^{-2}).
\]
The shift supplies the minus sign; the total twist is two.
The source computation is Figure~2, p.~34, and
\(\S5.1.2\), p.~29 \cite{VIII}.
It is not a newly computed full multiplication table.

A relative version retaining the actual boundary maps is
\[
 (\mathcal K\boxtimes\mathcal K)\otimes
                   R\Gamma(\mathbb P^2,E)(2)[1].
\]
The preceding nonconstancy dimension comparison is multiplied
by the three-dimensional total cohomology of \(\mathbb P^2\),
so it remains nonconstant. These two displayed coefficient
constructions use integral Tate twists; no geometric half
twist is needed.

The normalized positive factor \(v^2+1+v^{-2}\) here is
the proper Lefschetz object
\(R\Gamma(\mathbb P^2,E)(1)\).
Its three even cohomology groups and hyperplane maps give
the factor directly. This is a construction for the already
known printed coefficient. Equality of its trace with a
source constant does not identify the designated canonical
labels or prove a nonstandard structure-constant comparison.

\paragraph{Powers and the limit of tensor normalization.}
The product of \(a\) copies of the torsor gives an actual
\((\mathbb G_m)^a\)-torsor over \(O^a\), with coefficient
\(\mathcal K^{\boxtimes a}\).
Compactification by the product of its projective-line
bundles gives a cube of localization maps. Their totalization
and Koszul signs are the derived tensor product, so these
powers carry coherent ambient operations.

Let \(U\) be the arithmetic Weil line with geometric
Frobenius \(\sqrt q>0\) under the chosen complex embedding.
Then
\[
                   U^{-a}\otimes\mathcal K^{\boxtimes a}
\]
has fiber trace \((v-v^{-1})^a\) for \(v=q^{n/2}\).
For even \(a\) the normalization is the twist \((a/2)\).
For odd \(a\) it is an arithmetic line of odd integral weight,
with no geometric integral Tate object asserted.
An odd cohomological shift supplies an overall negative sign.

There is no tensor inverse of the torus coefficient in the
bounded finite-dimensional derived coefficient category.
Indeed the K\"unneth formula gives
\[
 d(A\otimes^L B)=d(A)d(B),\qquad
 d(A)=\sum_i\dim H^i(A).
\]
An invertible object has \(d=1\), so it is a single
one-dimensional line in one degree.
The total dimension of \(R\Gamma_c(\mathbb G_m,E)\) is two,
and that of its \(a\)-fold tensor power is \(2^a\).
Neither is invertible. A relative inverse for \(\mathcal K\)
would pull back to an impossible inverse at a geometric point.
The compact-support K\"unneth and \(\ell\)-adic passages are
\cite[Theorem~22.1 and Remark~22.3(b)--(d),
pp.~135--136]{Milne}.

Thus division by \((v-v^{-1})^a\) in a character ring is
not tensoring by an inverse of this finite geometric
coefficient object. An enlarged localization or a specified
extraction operation is additional data.

\paragraph{The exact comparison that remains.}
This test constructs standard ambient convolution maps and
the nonproper factors in two actual source coefficients.
It also proves an obstruction to the constant coefficient
replacement: the Chern-class boundary is nonzero although
all fiber traces agree with those of a constant object.

To recover canonical normalized positivity, one must still
identify the designated labels and their actual coefficient
objects, specify a relative factorization or a functorial
filtration/extraction compatible with these boundary and
composition maps, and identify the extracted factor with
proper pure Lefschetz data under the stated grading and
Frobenius comparison.
For the printed VIII coefficient, the projective-plane
factor supplies the last item; it does not supply the
canonical-label comparison.

The formal specialization condition also remains.
The two-torus coefficient has nonzero total cohomology
dimension four, even though its Euler character
\(c_0(1)\) is zero. Euler vanishing does not make the
object or its maps vanish. The source's cell refinement
requires its actual specialization or filtration comparison.

No nonstandard quadratic in the bounded ordinary Soergel
category has been constructed here, so \fid{F59} remains
unchanged. Nor is a full finite-algebra presentation by the
ordinary isomorphisms of \fid{F60} asserted. The valid
standard relation is a nonsplit triangle, and the compact
coefficient pushforwards are mixed. These scope limits are
part of the construction.

% Separate paragraph for Part III.
\paragraph{The nonproper kernel test in Grothendieck's language.}
The line bundle
\(\mathcal L=\operatorname{Hom}(L_x,L_z)\) and its
multiplicative-group torsor are actual geometric objects.
Their locally free sheaves use EGA~I, Theorem~1.3.7, p.~86,
Corollary~1.3.8, p.~87, and Theorem~1.4.1, pp.~90--92
\cite{EGA}. The first Chern class obstructs splitting of
the relative compact-pushforward triangle, as proved by
the two-section matrix \eqref{eq:twoSectionChernMatrix}.
This is actual extension data before passage to a
Grothendieck group.

Pullback, tensor product, compact pushforward, base change
and localization construct the displayed operation maps.
SGA~5, III, 4.7 supplies the trace-formula framework
\cite{SGA5}; equality of trace functions does not imply
an isomorphism of complexes, as the explicit dimension
comparison proves.

The coefficient gradings use SGA~3, I, 4.7.2--3 and
cote~90, scan~17; the effective/virtual distinction uses
SGA~6, V, \S\S2--3 and cote~133, scan~18
\cite{SGA3,SGA6,Archive}.
The factor \(v-v^{-1}\) is a virtual Euler character
of a nonproper compact pushforward, rather than an
effective graded dimension or an invertible finite
coefficient object. The bounded point complexes fit
SGA~6, I, Definition~4.7, pp.~121--122, and
Example~4.8, p.~122.
No \(\lambda\)-ring structure on the noncommutative
kernel algebra is asserted.

Pure coefficient sheaves on a torsor and their mixed
compact pushforwards are distinct objects. The proper
projective-plane factor has its usual Lefschetz
structure. None of these framework statements identifies
the constructions with the unknown nonstandard canonical
category.

